\chapter[Bericht des Kassirers der …]{Bericht des Kassirers der Allgemeinen Synode}
\quellseite{13}

vom 1. Februar 1903 bis 1. Februar 1904.

\section{Synodalkasse}

\subsection{Einnahme}

Bestand der Kasse am 1. Februar 1903: \$5376.03. Aus dem Illinois-District: \$3240.00. Aus dem Mittleren District: 3050.62. Aus dem Wisconsin-District: 1400.00. Aus dem Michigan-District: 1122.75. Aus dem Westlichen District: 1116.48. Aus dem Oestlichen District: 485.42. Aus dem Iowa-District: 259.12. Aus dem Kansas-District: 175.00. Aus dem Nebraska-District: 100.00. Aus dem Minnesota- und Dakota-District: 96.57. Aus dem Canada-District: 3.00. Schulgeld in Milwaukee: 380.00. Schulgeld in Hawthorne: 320.00. Schulgeld in St. Paul: 217.51. Verkauf eines Pferdes etc. in Addison: 7.50. Durch J. C. Koch von Dir. Albrecht in Milwaukee: 4.50. Seward, Nebr.: Gehalt zurückerstattet: 85.00. Ertrag eines Contos: 968.54. Concordia Publishing House: 71337.03. Summe: 84369.04. Total: \$89745.07.

\subsection{Ausgabe}

Gehalte der Herren Professoren und der Hausverwalter: in St. Louis \$9499.68, in Springfield 6224.68, in Addison 9499.44, in Milwaukee 8607.94, in Fort Wayne 9890.80, in St. Paul 5270.69, in Seward 2919.96, in Concordia 3600.00, in Hawthorne 3999.84. Extrabewilligung der Aufsichtsbehörde für Herrn Professor Bischoff: 100.00. Summe: \$59613.03.

\quellseite{14}

Uebertrag \$59613.03; \$89745.07.

Pension für Herrn Dr. H. C. Schwan 1000.00. Pension für Herrn Prof. L. Körner 450.00. Pension für Herrn Prof. G. \unsicher{Rohling} (auf Beschluß der Aufsichtsbehörde) 62.50. Pension für Herrn Prof. H. Wyneken 420.00. Pension für Herrn Prof. C. \unsicher{Frinckel} 420.00. Pension für Herrn Prof. C. Lange 420.00. Pension für Herrn P. A. Landmann 240.00. Reisekosten des Allgemeinen Präses Herrn Dr. F. Pieper 275.00. Reisekosten aller anderen Synodalbeamten 1248.38. Theil der Reisekosten der Synodalconferenz 6.00. Reisekosten der Revisionscommission 46.00. Umzugsunkosten Herrn Prof. J. M. F. Strieters nach Seward 233.77. Umzugsunkosten Herrn Prof. W. Wallers nach Hawthorne 15.00. Umzugsunkosten Herrn Professor G. Niemanns nach Seward 107.00.

Bibliothek in St. Louis 214.24. Bibliothek in Springfield (für 1903) 100.00. Bibliothek in Milwaukee (für 1903) 100.00. Bibliothek in Fort Wayne (für 1902 und 1903) 200.00. Bibliothek in Hawthorne (für 1903) 50.00. Bibliothek in St. Paul (für 1902 und 1903) 100.00. Alumnatskrankenfonds in Milwaukee (für 1903) 50.00. Alumnatskrankenfonds in Fort Wayne (für 1903) 50.00. Alumnatskrankenfonds in Addison (für 1903) 50.00. Alumnatskrankenfonds in St. Paul (für 1902 und 1903) 100.00. Alumnatskrankenfonds in Hawthorne (für 1903) 50.00. Unterstützung der Englischen Synode 3000.00. Deficit der New York-Emigrantenmissionskasse für 1903 666.00. Safe Deposit Box-Miethe 5.00. Rechnungen des Concordia Publishing House 9.50. Postage stamps, telegrams and express charges 17.65.

Von den Anstaltsbehörden verausgabt: in Fort Wayne 5223.17, in Milwaukee 3675.02, in St. Louis 3060.41, in Addison 2781.98, in Concordia 2323.34, in Springfield 2116.90, in St. Paul 1811.75, in Seward 975.33. 91306.97.

Bleibt Schuld am 1. Februar 1904 \$1561.90.

\section{Allgemeine Baukasse}

\subsection{Ausgabe}

Verbliebene Schuld am 1. Februar 1903 \$10726.97. Reparatur des Collegegebäudes und neue Heizung in Fort Wayne \$22707.73. Specialsteuer für Asphaltstraße an Maumee-Avenue in Fort Wayne 4817.29. Ein neues Gebäude in Concordia 4920.00. Einrichtung eines Gebäudes für Wohnzwecke in St. Paul 2073.03. Für einen Uebungssaal in St. Paul 1214.48. Specialsteuer für einen Uebungssaal, Wasser \&c. in Milwaukee 1270.50. Fire Escapes, neue Closets und \unsicher{Gaslicht-Ersetzung} in St. Louis 1168.45. Reparaturen in Springfield 855.25. Betten in Seward 665.20. 39722.33; \$50449.30.

\subsection{Einnahme}

Aus dem Illinois-District \$8670.12. Aus dem Mittleren District 6695.31. Aus dem Westlichen District 5017.57. Aus dem Wisconsin-District 4399.31. Aus dem Oestlichen District 4111.00. Aus dem Minnesota- und Dakota-District 3291.32. Aus dem Michigan-District 2880.67. Aus dem Nebraska-District 2064.79. Aus dem Iowa-District 1092.56. Aus dem Kansas-District 713.32. Aus dem Südlichen District 334.45. Aus dem Canada-District 252.48. Aus dem California- und Nevada-District 222.00. Aus dem Oregon- und Washington-District 62.25. Aus der Norwegischen Synode 5.00. 40688.95.

Bleibt Schuld am 1. Februar 1904 \$9760.35.

\section{Allgemeine Innere Mission}

\subsection{Einnahme}

Bestand der Kasse am 1. Februar 1903 \$5942.13. Aus dem Illinois-District \$4725.00. Aus dem Mittleren District 2500.00. Aus dem Michigan-District 117.26. Aus dem Wisconsin-District 68.00. Aus dem Oestlichen District 59.86. Aus dem Westlichen District 19.10. Aus dem Iowa-District 17.04. Ertrag des Legats: „Aus den Ersparnissen einer Pfarrfrau“ durch P. Lochner 21.60. 7527.86; \$13469.99.

\subsection{Ausgabe}

An den Oregon- und Washington-District \$5653.78. An den Südlichen District 3234.43. An den California- und Nevada-District 2655.96. An den Kansas-District 1100.00. An den Canada-District 600.00. Rechnung des Concordia Publishing House 20.40. 13264.57.

Bleibt Kassenbestand am 1. Februar 1904 \$205.42.

\section{Judenmission}

\subsection{Einnahme}

Bestand der Kasse am 1. Februar 1903 \$1517.09. Aus dem Illinois-District \$334.53. Aus dem Mittleren District 149.35. Aus dem Wisconsin-District 137.05. Aus dem Michigan-District 126.84. Aus dem Westlichen District 92.01. Aus dem Nebraska-District 76.58. Aus dem Iowa-District 65.85. Aus dem Canada-District 18.50. Aus dem Kansas-District 16.90. Aus dem California- und Nevada-District 8.50. Aus dem Oregon- und Washington-District 1.00. Aus der Norwegischen Synode 62.50. Aus der deutschen Freikirche 36.73. Aus der Englischen Synode 16.00. 1142.34; \$2659.43.

\subsection{Ausgabe}

An Herrn W. F. Weber, Kassirer der Commission 1500.00.

Bleibt Kassenbestand am 1. Februar 1904 \$1159.43.

\section{Taubstummenmission}

\subsection{Einnahme}

Bestand der Kasse am 1. Februar 1903 \$1707.34. Aus dem Illinois-District \$790.58. Aus dem Minnesota- und Dakota-District 523.49. Aus dem Mittleren District 519.61. Aus dem Westlichen District 388.90. Aus dem Nebraska-District 367.55. Aus dem Wisconsin-District 194.97. Aus dem Oestlichen District 193.49. Aus dem Michigan-District 171.88. Aus dem Iowa-District 109.76. Aus dem Kansas-District 89.97. Aus dem Canada-District 25.56. Aus dem Südlichen District 17.35. Aus dem California- und Nevada-District 12.00. Aus dem Oregon- und Washington-District 3.75. Aus der deutschen Freikirche 6.57. 3413.43; \$5120.77.

\subsection{Ausgabe}

Auf Anweisungen der Commission 3956.29.

Bleibt Kassenbestand am 1. Februar 1904 \$1164.48.

\section{Indianermission}

\subsection{Einnahme}

Bestand der Kasse am 1. Februar 1903 \$49.40. Aus dem Wisconsin-District \$1042.11. Aus dem Illinois-District 831.22. Aus dem Westlichen District 508.33. \$2381.66; \$49.40.

\quellseite{15}

Uebertrag \$2381.66; \$49.40. Aus dem Westlichen District 375.62. Aus dem Mittleren District 368.30. Aus dem Nebraska-District 290.41. Aus dem Michigan-District 241.65. Aus dem Minnesota- und Dakota-District 238.84. Aus dem Kansas-District 147.00. Aus dem Canada-District 117.55. Aus dem Iowa-District 98.39. Aus dem Südlichen District 44.50. Aus dem Oregon- und Washington-District 19.79. Aus dem California- und Nevada-District 13.50. Aus der deutschen Freikirche 18.92. 4356.32; \$4405.72.

\subsection{Ausgabe}

An Herrn P. C. Hanow \$1711.24. An Herrn P. H. Geßner 1412.64. An Herrn P. H. F. Knuth (für Arizona) 5.00. 3128.88.

Bleibt Kassenbestand am 1. Februar 1904 \$1276.84.

\section{Heidenmission}

\subsection{Einnahme}

Bestand der Kasse am 1. Februar 1903 \$20899.70. Aus dem Illinois-District \$1240.47. Aus dem Mittleren District 957.18. Aus dem Westlichen District 595.11. Aus dem Minnesota- und Dakota-District 589.57. Aus dem Oestlichen District 506.30. Aus dem Wisconsin-District 503.32. Aus dem Michigan-District 425.11. Aus dem Nebraska-District 392.49. Aus dem Kansas-District 233.12. Aus dem Iowa-District 216.23. Aus dem Canada-District 182.63. Aus dem Südlichen District 47.20. Aus dem California- und Nevada-District 26.50. Aus dem Oregon- und Washington-District 2.00. Aus der deutschen Freikirche 187.44. Aus der Englischen Synode 7.25. 6122.92; \$27022.62.

\subsection{Ausgabe}

An Herrn P. H. Leonhardt, Kassirer der Commission \$8511.00. Rechnung des Concordia Publishing House 200.78. 8711.78.

Bleibt Kassenbestand am 1. Februar 1904 \$18310.84.

\section{Mission in Brasilien}

\subsection{Einnahme}

Bestand der Kasse am 1. Februar 1903 \$4070.68. Aus dem Illinois-District \$1578.71. Aus dem Mittleren District 1202.21. Aus dem Westlichen District 533.64. Aus dem Nebraska-District 499.73. Aus dem Wisconsin-District 435.25. Aus dem Minnesota- und Dakota-District 433.82. Aus dem Michigan-District 324.16. Aus dem Oestlichen District 323.16. Aus dem Iowa-District 243.59. Aus dem Kansas-District 124.73. Aus dem Canada-District 27.35. Aus dem Südlichen District 13.00. Aus dem California- und Nevada-District 2.00. Aus der deutschen Freikirche 26.01. 5950.36; \$10021.04.

\subsection{Ausgabe}

Auf Anweisungen der Commission 4216.80.

Bleibt Kassenbestand am 1. Februar 1904 \$5804.24.

\section{Mission in London}

\subsection{Einnahme}

Bestand der Kasse am 1. Februar 1903 \$229.05. Aus dem Illinois-District \$624.92. Aus dem Westlichen District 333.53. Aus dem Michigan-District 291.61. \$1250.06; \$229.05.

Uebertrag \$1250.06; \$229.05. Aus dem Mittleren District 266.55. Aus dem Oestlichen District 132.30. Aus dem Minnesota- und Dakota-District 121.48. Aus dem Nebraska-District 117.22. Aus dem Wisconsin-District 83.60. Aus dem Kansas-District 44.88. Aus dem Iowa-District 28.00. Aus dem Südlichen District 5.00. Aus dem Canada-District 1.50. Aus der Englischen Synode 2.00. 2052.59; \$2281.64.

\subsection{Ausgabe}

Auf Anweisungen der Commission 1010.00.

Bleibt Kassenbestand am 1. Februar 1904 \$1271.64.

\section{Mission in Australien und Neuseeland}

\subsection{Einnahme}

Bestand der Kasse am 1. Februar 1903 \$195.23. Aus dem Illinois-District \$357.45. Aus dem Westlichen District 243.91. Aus dem Mittleren District 220.42. Aus dem Michigan-District 163.50. Aus dem Wisconsin-District 114.23. Aus dem Minnesota- und Dakota-District 103.07. Aus dem Nebraska-District 82.69. Aus dem Oestlichen District 77.36. Aus dem Kansas-District 57.64. Aus dem Iowa-District 41.67. Aus dem California- und Nevada-District 1.00. Aus dem Canada-District 1.00. 1464.24; \$1659.47.

\subsection{Ausgabe}

Auf Anweisungen der Commission 1025.00.

Bleibt Kassenbestand am 1. Februar 1904 \$634.47.

\section{Mission in Berlin}

\subsection{Ausgabe}

Verbliebene Schuld am 1. Februar 1903 \$410.02. Durch P. Keyl übermittelt für 1903 325.00. \$735.02.

\subsection{Einnahme}

Aus dem Illinois-District \$71.33. Aus dem Michigan-District 56.52. Aus dem Westlichen District 32.00. Aus dem Minnesota- und Dakota-District 25.55. Aus dem Mittleren District 20.59. Aus dem Iowa-District 14.35. Aus dem Kansas-District 9.17. Aus dem Wisconsin-District 7.00. Aus dem Nebraska-District 5.60. Aus dem Oestlichen District 3.00. 245.11.

Bleibt Schuld am 1. Februar 1904 \$488.91.

\section{Emigrantenmission in New York}

\subsection{Einnahme}

Durch die Districtskassirer \$644.69. Aus der deutschen Freikirche 3.81. Durch die Districtskassirer für das Pilgerhaus 30.55. \$679.05.

\subsection{Ausgabe}

An P. S. Keyl \$679.05.

\section{Emigrantenmission in Baltimore}

\subsection{Einnahme}

Durch die Districtskassirer \$216.95. Aus der deutschen Freikirche 1.27. \$218.22.

\subsection{Ausgabe}

An Herrn H. Schüller \$218.22.

\quellseite{16}

\section{Esthen- und Lettenmission}

\subsection{Einnahme}

Durch die Districtskassirer: \$395.46.

\subsection{Ausgabe}

An Herrn R. F. Hassenreffer, Kassirer der Commission: \$395.46.

\section{Sächsische Freikirche}

\subsection{Einnahme}

Durch die Districtskassirer: \$1039.48. Durch die Districtskassirer für Hartenstein: 274.86. Aus der Englischen Synode: 1.00. Summe: \$1315.34.

\subsection{Ausgabe}

Durch P. S. Keyl übermittelt: \$1315.34.

\section{Dänische Freikirche}

\subsection{Einnahme}

Durch die Districtskassirer: \$691.56.

\subsection{Ausgabe}

Durch P. S. Keyl übermittelt: \$691.56.

\section{Hermannsburger Freikirche}

\subsection{Einnahme}

Durch die Districtskassirer: \$44.24.

\subsection{Ausgabe}

Durch P. S. Keyl übermittelt: \$44.24.

\section{Mission in Hamburg}

\subsection{Einnahme}

Durch die Districtskassirer: \$75.67.

\subsection{Ausgabe}

Durch P. S. Keyl übermittelt: \$75.67.

\section{Allgemeine Kirchbaukasse}

\subsection{Einnahme}

Bestand der Kasse am 1. Februar 1903: \$1704.00. Aus dem Michigan-District: \$6.50. Aus dem Westlichen District: 6.00. Summe: 12.50. Gesamt: \$1716.50.

\subsection{Ausgabe}

Gemeinden geliehen: 1200.00. Bleibt Bestand der Kasse am 1. Februar 1904: \$516.50.

\section{Legatfonds der Synode}

\subsection{Einnahme}

Bestand der Kasse am 1. Februar 1903: \$4674.93. Johann Voßler-Legat, Auglaize Co., Ohio: 25.00. Johann D. Meyer-Legat, Crete, Ill.: 50.00. Maria Boye-Legat, St. Paul, Ill.: 800.00. Augusta Noack-Legat, St. Louis, Mo.: 50.00. Gesamt: \$5599.93.

\subsection{Ausgabe}

Keine. Bleibt Kassenbestand am 1. Februar 1904: \$5599.93.

\section{Legatfonds des Seminars zu St. Louis}

\subsection{Einnahme}

Bestand der Kasse am 1. Februar 1903: \$25452.30. Ertrag: 691.63. Gesamt: \$26143.93.

\subsection{Ausgabe}

Keine. Bleibt Kassenbestand am 1. Februar 1904: \$26143.93.

\section{Zusammenstellung}

Bestand der Kasse für Allgemeine Innere Mission: \$205.42. Bestand der Kasse für Judenmission: 1159.43. Bestand der Kasse für Taubstummenmission: 1164.48. Bestand der Kasse für Indianermission: 1276.84. Bestand der Kasse für Heidenmission: 18310.84. Bestand der Kasse für Mission in Brasilien: 5804.24. Bestand der Kasse für Mission in London: 1271.64. Bestand der Kasse für Australien und Neuseeland: 634.47. Bestand der Allgemeinen Kirchbaukasse: 516.50. Bestand des Legatfonds der Synode: 5599.93. Bestand des Legatfonds des Seminars zu St. Louis: 26143.93. Sonstige Gelder für verschiedene Zwecke bestimmt: 147.75. Gelder der Synode gegen Noten zinsfrei geliehen: 8720.00. Gesammt: \$70955.5\unsicher{…}.

Von diesen Geldern sind: Bar auf der Bank \$33005.32. In Bonds 18500.00. In Noten 1060.00. An zwei Gemeinden geliehen 4882.25. Einer Aufsichtsbehörde vorgestreckt 558.58. Der Australischen Synode vorgestreckt (Reisekosten Prof. C. F. Grähners) 1338.20. Summe: 59344.35. Bleibt Schuld am 1. Februar 1904: \$11611.16.

Die sich vertheilt wie folgt: Schuld in der Synodalkasse \$1561.90. Schuld in der Allgemeinen Baukasse 9560.35. Schuld in der Kasse für Mission in Berlin 488.91. Summe: \$11611.16.

St. Louis, Mo., den 1. Februar 1904.

J. F. Schuricht,

Kassirer der Allgemeinen Synode.
